% =============================================================================
%  Memoria del TFG — Backend de Sistema de Corrección Digital de Exámenes Escritos
%  Autor: Anxo Canay Reguera
%  Tutor: Eloy Anguiano Rey
%  Grado en Ingeniería Informática — EPS, Universidad Autónoma de Madrid
%  Curso: 2025/2026
% =============================================================================
%  Receta de compilación local (orden):
%    pdflatex main
%    bibtex   main
%    makeglossaries main
%    makeindex -s main.ist main
%    pdflatex main
%    pdflatex main
% =============================================================================

% arara: clean: {files: [main.aux, main.idx, main.ilg, main.ind, main.bbl, main.bcf, main.blg, main.run.xml, main.fdb_latexmk, main.fls, main.loe, main.lof, main.lol, main.lot, main.ltb, main.out, main.toc, main.upa, main.upb, main.acn, main.acr, main.alg, main.glg, main.glo, main.gls, main.glsdefs, main.xdy, main.loa, main.gnuploterrors, main.mw]}
% arara: pdflatex
% arara: makeglossaries
% arara: makeindex: {style: main.ist}
% arara: bibtex
% arara: pdflatex
% arara: pdflatex


\documentclass[epsbased,tfg,covers,final,printable,extendedindex,%
               firstnumbered,lof,lot,loc,loa]{tfgtfmthesisuam}


% =============================================================================
%  Preámbulo: metadatos, acrónimos y definiciones (modular)
% =============================================================================
% =============================================================================
%  Metadatos del documento
% =============================================================================
%  Variables del preámbulo de la plantilla tfgtfmthesisuam.
% =============================================================================

% --- Identificación académica ---------------------------------------------------
\title{Backend de Sistema de Corrección Digital de Exámenes Escritos}
\author{Anxo Canay Reguera}
\advisor{Eloy Anguiano Rey}
\levelin{Grado en Ingeniería Informática}
\copyrightdate{2026}


% --- Datos para portada y copyright -------------------------------------------
\coverdata{%
  Escuela Politécnica Superior\\
  Universidad Autónoma de Madrid\\
  C\textbackslash\ Francisco Tomás y Valiente nº 11%
}


% --- Ficheros de las secciones preliminares -----------------------------------
\resumenfile{front-matter/resumen}
\abstractfile{front-matter/abstract}

% --- Directorios de recursos --------------------------------------------------
\graphicsdir{files}
\logosdir{logos}
\codesdir{codes}
\datadir{data}


% --- Bibliografía -------------------------------------------------------------
%  Estilo: IEEE (bibtex 'ieeetr') · Fichero: bibliography/references.bib
\bibliographyconfig[ieeetr]{bibliography/references}


% --- Metadatos del PDF --------------------------------------------------------
\pdfmetavalues
  {Anxo Canay Reguera}
  {Backend de Sistema de Corrección Digital de Exámenes Escritos}%
  {Trabajo Fin de Grado}%
  {corrección digital, exámenes, OCR, Django, backend, sistema distribuido, multi-tenant}

% preamble/acronyms.tex
% Definición centralizada de todos los acrónimos del documento.
% Usar \ac{label} en primera aparición, \acs{label} para forma corta.
% ============================================================================

% --- Acrónimos del proyecto ---
\newacronym{scdee}{SCDEE}{Sistema de Corrección Digital de Exámenes Escritos}

% --- Protocolos y estándares web ---
\newacronym{api}{API}{Application Programming Interface}
\newacronym{rest}{REST}{Representational State Transfer}
\newacronym{http}{HTTP}{HyperText Transfer Protocol}
\newacronym{https}{HTTPS}{HyperText Transfer Protocol Secure}
\newacronym{cors}{CORS}{Cross-Origin Resource Sharing}
\newacronym{jwt}{JWT}{JSON Web Token}
\newacronym{tls}{TLS}{Transport Layer Security}
\newacronym{smtp}{SMTP}{Simple Mail Transfer Protocol}
\newacronym{sftp}{SFTP}{SSH File Transfer Protocol}
\newacronym{xss}{XSS}{Cross-Site Scripting}
\newacronym{crud}{CRUD}{Create, Read, Update, Delete}
\newacronym{sso}{SSO}{Single Sign-On}
\newacronym{oauth2}{OAuth2}{Open Authorization 2.0}
\newacronym{saml}{SAML}{Security Assertion Markup Language}
\newacronym{oidc}{OIDC}{OpenID Connect}

% --- Formatos de datos ---
\newacronym{json}{JSON}{JavaScript Object Notation}
\newacronym{csv}{CSV}{Comma-Separated Values}
\newacronym{pdf}{PDF}{Portable Document Format}
\newacronym{yaml}{YAML}{YAML Ain't Markup Language}
\newacronym{svg}{SVG}{Scalable Vector Graphics}
\newacronym{url}{URL}{Uniform Resource Locator}

% --- Tecnologías y frameworks ---
\newacronym{drf}{DRF}{Django REST Framework}
\newacronym{orm}{ORM}{Object-Relational Mapping}
\newacronym{ocr}{OCR}{Optical Character Recognition}
\newacronym{sql}{SQL}{Structured Query Language}
\newacronym{siem}{SIEM}{Security Information and Event Management}
\newacronym{asr}{ASR}{Automatic Speech Recognition}
\newacronym{ttl}{TTL}{Time To Live}

% --- Criptografía y seguridad ---
\newacronym{aes}{AES}{Advanced Encryption Standard}
\newacronym{gcm}{GCM}{Galois/Counter Mode}

% --- Identificadores ---
\newacronym{dni}{DNI}{Documento Nacional de Identidad}
\newacronym{nia}{NIA}{Número de Identificación del Alumno}

% --- Institucional ---
\newacronym{eps}{EPS}{Escuela Politécnica Superior}
\newacronym{uam}{UAM}{Universidad Autónoma de Madrid}
\newacronym{tfg}{TFG}{Trabajo Fin de Grado}

% --- Ingeniería del software ---
\newacronym{uml}{UML}{Unified Modeling Language}
\newacronym{er}{E-R}{Entidad-Relación}
\newacronym{rf}{RF}{Requisito Funcional}
\newacronym{rnf}{RNF}{Requisito No Funcional}

% --- Normativa ---
\newacronym{rgpd}{RGPD}{Reglamento General de Protección de Datos}
\newacronym{lopdgdd}{LOPDGDD}{Ley Orgánica de Protección de Datos y Garantía de Derechos Digitales}
\newacronym{arco}{ARCO}{Acceso, Rectificación, Cancelación y Oposición}

% --- Infraestructura ---
\newacronym{s3}{S3}{Simple Storage Service}
\newacronym{ci}{CI}{Continuous Integration}
\newacronym{cd}{CD}{Continuous Deployment}

% TODO: Añadir nuevos acrónimos aquí conforme aparezcan en el texto.
%       Formato: \newacronym{label}{SIGLA}{Significado Completo}

% =============================================================================
%  Definiciones
% =============================================================================

\newdefinition{backend}{backend}{backends}%
  {Conjunto de componentes de un sistema responsables de la lógica de negocio,
   la persistencia de datos y la exposición de servicios, sin contacto directo
   con la interfaz de usuario}

\newdefinition{frontend}{frontend}{frontends}%
  {Capa del sistema responsable de la interacción con el usuario y de la
   presentación visual de la información}

\newdefinition{multitenant}{multi-tenant}{multi-tenant}%
  {Arquitectura que permite atender a múltiples organizaciones (tenants)
   desde una única instancia del sistema, manteniendo aislados sus datos y
   configuraciones}

\newdefinition{plugin}{plugin}{plugins}%
  {Componente intercambiable que extiende la funcionalidad del sistema
   cumpliendo una interfaz definida, sin requerir modificaciones en el resto
   del código}

\newdefinition{tenant}{tenant}{tenants}%
  {Cada una de las organizaciones que comparten una instancia del sistema en
   una arquitectura multi-tenant}

\newdefinition{instancia}{instancia de examen}{instancias de examen}{Copia física
digitalizada de un examen asociada a un estudiante concreto, que incluye el
\acs{pdf} escaneado, sus anotaciones, calificaciones y metadatos}

\newdefinition{modelo}{modelo de examen}{modelos de examen}{Variante estructural
de un examen que define el número de problemas, sus puntuaciones, las zonas de
reconocimiento y el \acs{pdf} en blanco de referencia}

\newdefinition{anotacion}{anotación}{anotaciones}{Registro creado por un corrector
sobre una \dfn{instancia} que puede contener texto enriquecido, trazos de
\textit{stylus} o audio, y que opcionalmente se asocia a un problema concreto}

\newdefinition{rubrica}{rúbrica}{rúbricas}{Conjunto de criterios de evaluación con
puntuación fija asociados a un problema, cuyos \textit{checkmarks} el corrector
marca para calcular automáticamente la calificación}

\newdefinition{corrector}{corrector}{correctores}{Profesor al que se le han asignado
una o más \dfnpl{instancia} o problemas de un examen para su evaluación}

\newdefinition{coordinador}{coordinador}{coordinadores}{Profesor con el rol principal
de gestión dentro de una asignatura, responsable de crear exámenes, asignar
correctores, gestionar revisiones y publicar calificaciones}

\newdefinition{gestor}{gestor}{gestores}{Usuario con permisos de administración de la
plataforma, responsable de la gestión de usuarios, cursos académicos y configuración
global del sistema}

\newdefinition{watcher}{watcher}{watchers}{Proceso demonio que monitoriza una o varias
carpetas \acs{sftp} en busca de nuevos ficheros \acs{pdf} para iniciar su ingesta
automática}

\newdefinition{softdelete}{desactivación lógica}{desactivaciones lógicas}{Mecanismo de
baja de un registro que preserva su existencia en la base de datos mediante un campo
indicador, garantizando la integridad referencial de los datos históricos}

\newdefinition{zonaocr}{zona de reconocimiento}{zonas de reconocimiento}{Región
rectangular definida por coordenadas y página sobre el \acs{pdf} en blanco de un
\dfn{modelo}, que indica al motor \acs{ocr} dónde buscar un atributo concreto}

\newdefinition{regla}{regla de asignación}{reglas de asignación}{Configuración que
define qué \dfnpl{corrector} evalúan qué \dfnpl{instancia} o problemas, basándose en
filtros como grupo, \dfn{modelo} o problema}

\newdefinition{endpoint}{endpoint}{endpoints}{Punto de acceso de una \ac{api} \ac{rest},
identificado por una \ac{http} y una ruta, que acepta peticiones y devuelve respuestas}

\newdefinition{worker}{worker}{workers}{Proceso independiente que ejecuta tareas
asíncronas extraídas de una cola de mensajes}

\newdefinition{broker}{broker}{brokers}{Intermediario de mensajes que gestiona la
comunicación entre productores y consumidores de tareas asíncronas}

\newdefinition{pipeline}{pipeline}{pipelines}{Secuencia encadenada de etapas de
procesamiento en la que la salida de una etapa alimenta la entrada de la siguiente}

\newdefinition{bucket}{bucket}{buckets}{Contenedor lógico de almacenamiento de objetos en un
servicio de almacenamiento, utilizado para organizar y gestionar los archivos del sistema}

\newdefinition{multi-tenant}{multi-tenant}{multi-tenant}{Arquitectura de software en la que una única
instancia de la aplicación sirve a múltiples organizaciones o clientes, manteniendo la
separación de datos y configuraciones entre ellos}

\newdefinition{single-tenant}{single-tenant}{single-tenant}{Arquitectura de software en la que cada
organización o cliente tiene su propia instancia dedicada de la aplicación, con su propia
base de datos y configuración}

\newdefinition{polling}{polling}{polling}{Técnica de comunicación en la que un proceso cliente
realiza peticiones periódicas a un servidor para comprobar si hay nuevos datos o eventos disponibles}



% =============================================================================
%  Documento
% =============================================================================
\begin{document}


% -----------------------------------------------------------------------------
%  Capítulo 1 — Introducción
% -----------------------------------------------------------------------------
\chapter{Introducción\label{CAP:INTRODUCTION}}
  \section{Motivación\label{SEC:INTRO-MOTIVATION}}%
          {chapters/01-introduction/11-motivation}
  \section{Objetivos\label{SEC:INTRO-OBJECTIVES}}%
          {chapters/01-introduction/12-objectives}
  \section[Contexto del proyecto y coordinación]%
          {Contexto del proyecto y coordinación con el \acs{tfg} de
           \textit{frontend}\label{SEC:INTRO-CONTEXT}}%
          {chapters/01-introduction/13-context}
  \section{Estructura del documento\label{SEC:INTRO-STRUCTURE}}%
          {chapters/01-introduction/14-structure}


% -----------------------------------------------------------------------------
%  Capítulo 2 — Estado del arte
% -----------------------------------------------------------------------------
\chapter{Estado del arte\label{CAP:STATE-OF-THE-ART}}
  \section[Plataformas de corrección digital]%
          {Plataformas de corrección digital de
           exámenes\label{SEC:SOTA-PLATFORMS}}%
          {chapters/02-state-of-the-art/21-platforms}
    \subsection{Crowdmark\label{SS:SOTA-CROWDMARK}}%
                {chapters/02-state-of-the-art/211-crowdmark}
    \subsection{Gradescope\label{SS:SOTA-GRADESCOPE}}%
                {chapters/02-state-of-the-art/212-gradescope}
    \subsection{Akindi\label{SS:SOTA-AKINDI}}%
                {chapters/02-state-of-the-art/213-akindi}
    \subsection[Integración con los LMS]%
                {El papel central de la integración con los \emph{Learning
                 Management Systems}\label{SS:SOTA-LMS}}%
                {chapters/02-state-of-the-art/214-lms}
    \subsection[Corrección asistida por modelos de lenguaje]%
                {La ola de la corrección asistida por modelos de
                 lenguaje\label{SS:SOTA-LLM}}%
                {chapters/02-state-of-the-art/215-llm}
    \subsection{Posicionamiento del SCDEE\label{SS:SOTA-POSITIONING}}%
                {chapters/02-state-of-the-art/216-positioning}
  \section{Tecnologías base\label{SEC:SOTA-TECH}}%
          {chapters/02-state-of-the-art/22-technologies}
    \subsection[Django y DRF]%
                {Django y Django {REST} Framework\label{SS:SOTA-DJANGO}}%
                {chapters/02-state-of-the-art/221-django}
    \subsection{PostgreSQL\label{SS:SOTA-POSTGRES}}%
                {chapters/02-state-of-the-art/222-postgresql}
    \subsection{Celery y Redis\label{SS:SOTA-CELERY}}%
                {chapters/02-state-of-the-art/223-celery}
    \subsection{MinIO\label{SS:SOTA-MINIO}}%
                {chapters/02-state-of-the-art/224-minio}
    \subsection{Tesseract y EasyOCR\label{SS:SOTA-OCR}}%
                {chapters/02-state-of-the-art/225-ocr}
  \section{Conceptos transversales\label{SEC:SOTA-CROSS}}%
          {chapters/02-state-of-the-art/23-cross-concepts}
    \subsection[Patrones multi-tenant]%
                {Patrones \textit{multi-tenant}\label{SS:SOTA-MULTITENANT}}%
                {chapters/02-state-of-the-art/231-multitenancy}
    \subsection[Arquitecturas basadas en plugins]%
                {Arquitecturas basadas en
                 \emph{plugins}\label{SS:SOTA-PLUGINS}}%
                {chapters/02-state-of-the-art/232-plugins}


% -----------------------------------------------------------------------------
%  Capítulo 3 — Análisis
% -----------------------------------------------------------------------------
\chapter{Análisis\label{CAP:ANALYSIS}}
  \section{Proceso de análisis y actores\label{SEC:ANALYSIS-PROCESS}}%
          {chapters/03-analysis/31-process}
  \section{Requisitos funcionales\label{SEC:ANALYSIS-RF}}%
          {chapters/03-analysis/32-functional}
  \section{Requisitos no funcionales\label{SEC:ANALYSIS-RNF}}%
          {chapters/03-analysis/33-nonfunctional}


% -----------------------------------------------------------------------------
%  Capítulo 4 — Diseño
% -----------------------------------------------------------------------------
\chapter{Diseño\label{CAP:DESIGN}}{chapters/04-design/intro}
  \section{Arquitectura general\label{SEC:DES-OVERVIEW}}%
          {chapters/04-design/41-overview}
    \subsection{Componentes\label{SS:DES-COMPONENTS}}%
                {chapters/04-design/411-components}
    \subsection{Modelo de despliegue\label{SS:DES-DEPLOYMENT}}%
                {chapters/04-design/412-deployment}
  \section[Modelo de datos y multi-tenant]%
          {Modelo de datos y aislamiento
           \dfn{multitenant}\label{SEC:DES-DATA}}%
          {chapters/04-design/42-data}
    \subsection{Estructura del modelo\label{SS:DES-MODEL-STRUCTURE}}%
                {chapters/04-design/421-structure}
    \subsection{Visión global del modelo\label{SS:DES-MODEL-OVERVIEW}}%
                {chapters/04-design/422-overview}
    \subsection[Aislamiento multi-tenant]%
                {Aislamiento \dfn{multitenant}%
                 \label{SS:DES-MODEL-ISOLATION}}%
                {chapters/04-design/423-isolation}
  \section[Capa de servicios y arquitectura modular]%
          {Capa de servicios y arquitectura
           modular\label{SEC:DES-LAYERS}}%
          {chapters/04-design/43-layers}
    \subsection[Arquitectura por capas]%
                {Arquitectura por capas con
                 \textit{Service Layer}\label{SS:DES-SERVICE-LAYER}}%
                {chapters/04-design/431-service-layer}
    \subsection[Aplicaciones Django]%
                {Organización modular en aplicaciones
                 Django\label{SS:DES-APPS}}%
                {chapters/04-design/432-apps}
    \subsection[Plugins intercambiables]%
                {\dfnpl{plugin} intercambiables\label{SS:DES-PLUGINS}}%
                {chapters/04-design/433-plugins}
  \section[Pipeline de ingesta y reconocimiento]%
          {\textit{Pipeline} de ingesta y
           reconocimiento\label{SEC:DES-INGEST}}%
          {chapters/04-design/44-pipeline}
    \subsection{Detección y encolado\label{SS:DES-INGEST-DETECTION}}%
                {chapters/04-design/441-detection}
    \subsection{Procesamiento\label{SS:DES-INGEST-PROCESSING}}%
                {chapters/04-design/442-processing}
    \subsection{Reconocimiento\label{SS:DES-INGEST-RECOGNITION}}%
                {chapters/04-design/443-recognition}
    \subsection{Asignación\label{SS:DES-INGEST-ASSIGNMENT}}%
                {chapters/04-design/444-assignment}
  \section{Ciclo de vida de las instancias\label{SEC:DES-LIFECYCLE}}%
          {chapters/04-design/45-lifecycle}
  \section{Operaciones del sistema\label{SEC:DES-OPS}}%
          {chapters/04-design/46-operations}
    \subsection{Configuración del examen\label{SS:DES-OPS-EXAM}}%
                {chapters/04-design/461-exam-config}
    \subsection[Convocatoria y asignación de correctores]%
                {Convocatoria y asignación de
                 correctores\label{SS:DES-OPS-CONVOCATION}}%
                {chapters/04-design/462-convocation-correctors}
    \subsection{Anotaciones y calificación\label{SS:DES-OPS-ANNOT}}%
                {chapters/04-design/463-annotations-grading}
    \subsection{Operaciones de organización\label{SS:DES-OPS-ORG}}%
                {chapters/04-design/464-organization-ops}
  \section{Seguridad y cumplimiento\label{SEC:DES-SECURITY}}%
          {chapters/04-design/47-security}
  \section[API REST y contrato de integración]%
          {\acs{api} \acs{rest} y contrato de
           integración\label{SEC:DES-API}}%
          {chapters/04-design/48-api}


% -----------------------------------------------------------------------------
%  Capítulos 5-7 — Implementación, Pruebas y Conclusiones (placeholders)
% -----------------------------------------------------------------------------
\chapter{Implementación\label{CAP:IMPLEMENTATION}}
  \section{Metodología y plan de trabajo\label{SEC:IMPL-METHODOLOGY}}%
          {chapters/05-implementation/51-methodology}
    \subsection{Metodología seguida\label{SS:IMPL-METHOD}}%
                {chapters/05-implementation/511-method}
    \subsection{Plan de trabajo\label{SS:IMPL-PLAN}}%
                {chapters/05-implementation/512-plan}
    \subsection{Cronograma\label{SS:IMPL-GANTT}}%
                {chapters/05-implementation/513-gantt}
  \section{Entorno y herramientas\label{SEC:IMPL-TOOLING}}%
          {chapters/05-implementation/52-tooling}
  \section[Aspectos críticos de implementación]%
          {Aspectos críticos de
           implementación\label{SEC:IMPL-CRITICAL}}%
          {chapters/05-implementation/53-critical}
  \section{Despliegue\label{SEC:IMPL-DEPLOYMENT}}%
          {chapters/05-implementation/54-deployment}
\chapter{Pruebas y validación\label{CAP:TESTING}}%
        {chapters/06-testing/intro}
  \section{Estructura de pruebas\label{SEC:TEST-STRUCTURE}}%
          {chapters/06-testing/61-test-structure}
  \section{Cobertura\label{SEC:TEST-COVERAGE}}%
          {chapters/06-testing/62-coverage}
  \section{Pruebas de carga\label{SEC:TEST-LOAD}}%
          {chapters/06-testing/63-load-tests}
  \section[Calidad del reconocimiento óptico]%
          {Calidad del reconocimiento
           óptico\label{SEC:TEST-OCR-QUALITY}}%
          {chapters/06-testing/64-ocr-quality}
  \section{Validación funcional\label{SEC:TEST-FUNCTIONAL}}%
          {chapters/06-testing/65-functional}
\chapter{Conclusiones y trabajo futuro\label{CAP:CONCLUSIONS}}
  \section{Conclusiones\label{SEC:CONCL-CONCLUSIONS}}%
          {chapters/07-conclusions/71-conclusions}
  \section{Trabajo futuro\label{SEC:CONCL-FUTURE-WORK}}%
          {chapters/07-conclusions/72-future-work}


% -----------------------------------------------------------------------------
%  Apéndices
% -----------------------------------------------------------------------------
\appendix

\chapter{Requisitos del sistema\label{AP:REQUIREMENTS}}%
        {appendices/A-system-requirements/A-overview}
  \section{Requisitos funcionales\label{SEC:REQ-FUNCTIONAL}}%
          {appendices/A-system-requirements/A1-functional}
  \section{Requisitos no funcionales\label{SEC:REQ-NONFUNCTIONAL}}%
          {appendices/A-system-requirements/A2-nonfunctional}

\chapter{Manual de instalación y despliegue\label{AP:INSTALLATION}}%
        {appendices/B-installation/B-overview}
  \section{Prerrequisitos\label{SEC:INSTALL-PRE}}%
          {appendices/B-installation/B1-prerequisites}
  \section{Configuración inicial\label{SEC:INSTALL-CONFIG}}%
          {appendices/B-installation/B2-configuration}
  \section{Primera ejecución\label{SEC:INSTALL-FIRST-RUN}}%
          {appendices/B-installation/B3-first-run}
  \section{Operaciones habituales\label{SEC:INSTALL-OPS}}%
          {appendices/B-installation/B4-common-ops}
  \section{Acceso a los servicios\label{SEC:INSTALL-SERVICES}}%
          {appendices/B-installation/B5-services-access}
  \section{\textit{Reset} del entorno\label{SEC:INSTALL-RESET}}%
          {appendices/B-installation/B6-reset}

\chapter{Resultados detallados de pruebas\label{AP:TEST-RESULTS}}%
        {appendices/C-detailed-test-results/C-overview}
  \section{Pruebas de carga\label{SEC:CDETAIL-LOAD}}%
          {appendices/C-detailed-test-results/C1-load-tests}
  \section[Cobertura por aplicación]%
          {Cobertura por aplicación\label{SEC:CDETAIL-COVERAGE}}%
          {appendices/C-detailed-test-results/C2-coverage}
  \section[Calidad del reconocimiento óptico]%
          {Calidad del reconocimiento óptico\label{SEC:CDETAIL-OCR}}%
          {appendices/C-detailed-test-results/C3-ocr-quality}

\chapter{Catálogo de servicios por aplicación\label{AP:SERVICES}}%
        {appendices/D-services/D-overview}

\chapter{Especificación OpenAPI\label{AP:OPENAPI}}%
        {appendices/E-openapi/E-overview}
  \section{Catálogo de recursos\label{SEC:OPENAPI-CATALOG}}%
          {appendices/E-openapi/E1-catalog}
  \section{Acceso a la especificación\label{SEC:OPENAPI-ACCESS}}%
          {appendices/E-openapi/E2-access}

\chapter{Diagramas de clase completos\label{AP:CLASS-DIAGRAMS}}%
        {appendices/F-class-diagrams/F-overview}
  \section{Diagramas por subdominio\label{SEC:CLASS-BY-DOMAIN}}%
          {appendices/F-class-diagrams/F1-by-domain}
  \section{Vista global\label{SEC:CLASS-GLOBAL}}%
          {appendices/F-class-diagrams/F2-global}


\end{document}
